\chapter{讨论}
\label{ch:discussion}

\section{对假设的评估}
\label{sec:disc-hyp}

\paragraph{H1（降阶优化）}在模型范围内得到支持。只要模型中包含升力且未被压制，全部101个可行最优解都是升力推进的（lift-propelled）。H1不仅在模型之间得到检验，即升力推进步态（gait）所需功率为最佳纯阻力步态的二分之一到三分之一；也在同一模型内得到检验：以升力占比限值强制的阻力型划动需要四至十倍的功率，且达不到\SI{0.30}{\metre\per\second}，不过这第二个倍数仅基于3个可行划动。如下文所述，这一结论以一个未经标定的升力模型为前提。

\paragraph{H2（阻力与升力随速度的对比）}得到支持。正如预期，经过优化的桨（paddle）在低速下明显更强，其推力（thrust）随速度急剧下降，在\SI{0.30}{\metre\per\second}时降至约零；尾鳍（fluke）则保留了更大比例的推力，并在\SIrange{0.23}{0.24}{\metre\per\second}以上变得更强。桨的速度极限的成因已经明确：一是回程（recovery）阻力，它从静止到\SI{0.30}{\metre\per\second}增大到原来的七倍以上；二是在迎面来流中开扇所需的动量；与此同时，划水冲程（power stroke）推力的下降不到一半。效率的交叉点（crossover）略低于推力的交叉点，位于约\SI{0.22}{\metre\per\second}；在此以下，经过优化的桨比尾鳍更强，效率也更高，在\SI{0.15}{\metre\per\second}时达到$\etaF=0.35$--0.37。这与生物学中的图景不同：在动物中，升力型游动者在所有速度下都比划水者更高效\cite{Fish1996}；而对于配备折叠性能良好的桨的小型机器人腿，这一点仅在交叉点以上成立。H2的第二部分得到证实：调度尾鳍的俯仰角（pitch），使攻角（angle of attack）保持在\SI{20}{\degree}附近，可将\SI{0.30}{\metre\per\second}时的效率从0.24提高到0.40。

桨的效率超过了Fish针对划水动物所报告的约0.33的最大值\cite{Fish1984,Fish1996}。这两个数值不能直接比较：动物的数值是针对整个动物的估计值，而此处的$\etaF$是一个孤立推进器（propulsor）的水动力弗劳德效率（Froude efficiency），该推进器收拢后的扇面所产生的力仅约为张开扇面的九分之一（$a=0.11$，而麝鼠的面积比为0.445）。它也仍低于\SI{0.15}{\metre\per\second}时的准定常上限$\U/v_p\approx0.6$，其中扇面形心的平均划水冲程速度为$v_p\approx\SI{0.25}{\metre\per\second}$。

\paragraph{H3（整机）}基本得到支持。采用阻力（resistance）的最佳估计值时，尾鳍达到\SI{0.30}{\metre\per\second}，桨达到\SIrange{0.26}{0.27}{\metre\per\second}，因此两者速度相近但不相等；而随航速调度俯仰的尾鳍所需功率比桨低43\,\%（\SI{42}{}对\SI{73}{\milli\watt}）。不过，首次硬件测试比预测慢1.3至1.6倍（\cref{sec:res-hw}）。因此，预测速度是现有机器人的上限，而推进器之间的比较以CFD为依据。

\section{降阶模型与CFD为何在低速下存在差异}
\label{sec:disc-models}

降阶模型（reduced-order model）与CFD一致表明：阻力型划水（drag-based paddling）存在速度极限，而升力型拍动（lift-based flapping）在巡航速度（cruising speed）下占优。两者在低速下存在分歧：MuJoCo发现升力推进的步态在\SI{0.10}{\metre\per\second}时就已更为经济，差距约为二至四倍（在\SI{0.20}{\metre\per\second}时可达十倍），而CFD则发现经过优化的桨在约\SI{0.22}{\metre\per\second}以下效率更高。以下四点差异可以解释这一现象。
\begin{enumerate}
  \item \emph{升力模型偏乐观。}平板律$c_l\sin2\alpha$在产生升力时不含有限翼的诱导阻力和翼尖损失，而 toy\_quad 模型中脚蹼的展弦比（aspect ratio）仅为1.4。确认（validation）算例提示，即使在展弦比为6时，有限展长也很可能是高达18\,\%的效率亏损的原因之一，翼尖损失估计为10--20\,\%。该模型也不包含动态失速、环量建立的延迟以及脚蹼与其自身尾流的相互作用。
  \item \emph{对于起动中的桨，阻力模型偏悲观。}恒定的$C_D=2.2$忽略了起动涡（starting vortex）和附加质量（added mass）效应，而在CFD中，这些效应将桨划水冲程的有效力系数提高到其准定常值的2--3倍（\cref{sec:res-drag}）。
  \item \emph{阻力型冲程不同。}MuJoCo中的脚蹼只能通过转动来顺桨（feathering），既不能折叠（fold），也没有受控的折叠时机，并且仅限于式\eqref{eq:gait}的平滑谐波运动。由于斜着划过水的脚蹼会产生升力，升力模型中的阻力型划动实际上是被压制了升力的划动，而纯阻力模型中的大多数划水型划动在加入升力后都变为升力推进（\cref{sec:mj-liftshare}）。相比之下，CFD中的桨是专门的阻力装置：它以板面正对运动方向运动，在经过优化的相位下折叠至12\,mm，并采用梯形速度规律，上述每项特征各自可带来最多约30\,\%的推力提升。因此，MuJoCo比较的是升力型划动与被压制升力的脚蹼划动，而CFD比较的是尾鳍与经过优化的折叠桨。
  \item \emph{优化投入在两个方向上都不对称。}在MuJoCo中，升力步态针对给定所需速度下的功率进行了优化，而纯阻力前沿尚未收敛，升力模型内针对阻力型划动的带约束搜索在30次运行中也只有3次找到可行解。在CFD中，桨经过了多轮优化，而尾鳍运动L0取自另一个模型，之后仅对其俯仰幅值进行了随航速的调度。两者的面积也不同（\SI{18.5}{}对\SI{11.2}{\square\centi\metre}）。
\end{enumerate}
因此，降阶模型最好理解为关于\emph{机理}的论据，即优化器在条件允许时会利用哪种推进原理；而CFD则是关于\emph{量级}以及一种原理超越另一种原理时所处速度的论据。在两者存在分歧之处，以解析了流场且经过确认的CFD为准。以CFD为基准标定降阶模型的系数，尤其是带展弦比修正的$c_l$，将使两者协调一致，这也是顺理成章的下一步。

\section{对BODY2的启示}
\label{sec:disc-design}
\label{sec:disc-limits-prop}

\paragraph{建议。}对于在水面横渡开阔水域的任务，推荐BODY2采用俯仰角随航速减小的尾鳍作为推进器。在最佳估计阻力下，它比经过优化的桨V3trap略快，所需水动力功率（hydrodynamic power）低43\,\%，并且既不需要受控折叠，也不需要梯形速度规律。与机器人目前的配置相比，它使预测的巡航速度提高到约三倍；在硬件上一次未计时的对比中，尾鳍也明显比桨游得快（\cref{sec:plat-hw}）。在加速、低速定点保持或行走较为重要的场合，桨仍是更好的选择；此时需要一种收拢宽度（folded width）减小至约\SI{12}{\milli\metre}、并能在特定相位于\SIrange{45}{60}{\milli\second}内受控完成的折叠，而目前的拉线折叠机构无法做到这一点。这种分工，即划桨用于起步和低速、拍动用于巡航，正是在动物中发现的分工\cite{Walker2000,Fish1996}；一条能够从桨切换为尾鳍的腿（如Baines等人\cite{Baines2022}的可变形龟类肢体）将能够覆盖这两种工况。

\paragraph{尾鳍。}尾鳍在静止时的推力有限，因为它处于失速状态：没有前进速度时，拍动翼（foil）就相当于一个小桨。有航速时，其性能取决于能否使攻角保持在高效范围附近，而固定俯仰规律（fixed pitch schedule）做不到这一点。调度规律$\theta_0=\gamma_\text{max}/2$是一个能够实现这一点的单参数规则，同时还消除了不匹配的规律所产生的较大平均垂向力。在硬件上，俯仰是被动的：它由绕销轴（pin）的水动力力矩（使尾鳍转向来流方向）与板簧（leaf spring）之间的平衡决定。因此，被动尾鳍能够适应航速，但适应的比例并不符合要求。水动力力矩随相对速度的平方增长，而弹簧力矩则不然，因此比值$\theta/\gamma$随航速增大。线性弹簧只能在一个设计航速下满足$\theta_0=\gamma_\text{max}/2$；低于该航速时，尾鳍俯仰不足，攻角过大；高于该航速时，尾鳍与来流过度对齐，攻角降至高效范围以下。要在一段航速范围内实现匹配，需要渐进刚度弹簧、可调预紧力或主动俯仰执行器。硬件测试仅达到CFD推力的大约四分之一到一半（\cref{sec:res-hw}），这与被动俯仰未调至高效攻角相符，但不能排除其他原因。

\paragraph{等面积。}本文所用尾鳍的面积为桨的61\,\%。静止时两者的单位面积推力相同（\SIrange{3.9}{4.5}{}对\SI{4.1}{\milli\newton\per\square\centi\metre}），而在\SI{0.223}{\metre\per\second}时尾鳍已产生更大的单位面积推力（\SI{1.9}{}对\SIrange{1.3}{1.4}{\milli\newton\per\square\centi\metre}）。若推力系数和功率系数保持不变，一个与桨面积相同（\SI{18.5}{\square\centi\metre}）、按相同运动学运动的尾鳍，在静止时将产生约\SI{75}{\milli\newton}的推力，与桨相当，并且从约\SI{0.15}{\metre\per\second}起超过桨。其功率将按相同比例增长，静止时约为\SI{30}{\milli\watt}，因此效率交叉点不会移动。已为该机器人设计了一个具有此面积的月牙形尾鳍（展长\SI{86}{\milli\metre}，$\mathit{AR}\approx4$）；其较大的展弦比应能减小翼尖损失，但其较大的铰链力矩（在相同俯仰规律下约为A2的1.6倍）必须由俯仰弹簧或执行器承担。这是估算值，而非仿真结果。

\paragraph{执行器。}所有运动的水动力力矩与舵机（servo）相比都很小：其中最大的是梯形桨冲程绕髋（hip）关节的\SI{36}{\milli\newton\metre}，为PTK\,7465W堵转力矩（stall torque，\SIrange{0.44}{0.57}{\newton\metre}）的6--8\,\%。真正起作用的限制是关节速度以及切换折叠的能力（\cref{tab:actuators}）。L0的给定俯仰需要\SI{712}{\degree\per\second}，超过\SI{6}{\volt}时的空载转速（no-load speed，\SI{650}{\degree\per\second}），并接近\SI{8.4}{\volt}时的空载转速（\SI{860}{\degree\per\second}）；而\SI{0.30}{\metre\per\second}时随航速调度的俯仰（speed-scheduled pitch）仅需\SI{366}{\degree\per\second}。

\section{局限性}
\label{sec:disc-limits}

\begin{itemize}
  \item \emph{单腿。}CFD仅模拟一个推进器，不含连杆、船体（hull）和其他腿。腿间相互作用可使四条划水腿的推力高于单腿值的四倍\cite{Wang2025}，且这种作用在桨与尾鳍之间可能有所不同。自航（self-propulsion）平衡忽略了这一点。
  \item \emph{自由液面。}所有针对推进器的CFD均为带刚性盖（rigid lid）的单相计算。在硬件上和Flare Live中，尾鳍在冲程顶端会冲出水面约\SI{7}{\milli\metre}，而桨的回程路径也接近水面；通气和波浪均未建模。
  \item \emph{给定的俯仰与折叠。}尾鳍俯仰以刚性方式给定，而硬件上的尾鳍是被动的。扇面的折叠被建模为两个分别模拟的状态之间的瞬时切换。切换时附加质量的变化在一定界限内通过解析方法修正，所用附加质量由CFD辨识得到；张开状态算例自身的回程尾流对其划水冲程的影响未作修正，这使桨的结果略偏乐观。
  \item \emph{CFD精度。}流动按层流处理，这是$\mathit{Re}=4\times10^4$的确认算例中出现18\,\%效率亏损的原因之一；在尾鳍所处的较低雷诺数下，预计该影响较小。尾鳍推力约有$\pm12$\,\%的离散不确定性。桨的网格更粗，且未单独进行研究；\cref{sec:ver-drag}中已论证，结论对该量级的误差不敏感。
  \item \emph{功率。}此处的功率是推进器的水动力输入功率（input power）。舵机损耗、腿的惯性和连杆的阻力均未计入，因此推进器之间的功率比并不等于电池功率之比。
  \item \emph{推进器不对等。}桨的冲程比尾鳍运动优化得更充分，因此尾鳍的结果偏保守，对沉浮（heave）幅值、频率和俯仰进行系统优化很可能使交叉点移向更低速度。两种推进器还以不同频率运行（\SI{1.39}{}和\SI{2}{\hertz}），面积也不同；文中给出了单位面积推力和单位功率推力，但未进行等面积、等频率的比较。
  \item \emph{阻力。}船体阻力仅在一定界限内已知，巡航速度也相应存在不确定性。$D_\text{mid}$中的连杆项是工程估算值，而非测量值；此外，船体阻力是按尖端在前计算的，而硬件以钝端在前游动。
  \item \emph{降阶模型。}MuJoCo模型是一个尺寸相近的替代模型，并非BODY2，且其系数未经标定。其舵机模型（\SI{3}{\newton\metre}，\SI{400}{\degree\per\second}）比PTK\,7465W（\SIrange{0.44}{0.57}{\newton\metre}，\SIrange{650}{860}{\degree\per\second}）力矩更大、速度更慢。其搜索未完全收敛，模型内比较仅基于30次搜索（两个种子）中的3个可行阻力型划动，因此其功率倍数可能被高估。
  \item \emph{硬件。}硬件测试仅包括每个频率一次从静止起步的手动计时，船体反装，俯仰为被动，冲程幅值未测量。它表明自航估计是上限，但无法确定差异的原因。
\end{itemize}
